\manualchapter{Control reference}{sec:controls}
The five waveform editors run down the left edge. Beside them are global
Active Channels and Waveform Morph controls, each with a depth knob and CV
input. Eight voice columns to the right contain frequency, bipolar FM depth,
level, corresponding CV inputs, and an audio output. Morph is shared by
all eight oscillators within each polyphonic chip instance.

\begin{tabularx}{\textwidth}{@{}l l X@{}}
\toprule Control & Range; default & CV and behavior \\\midrule
Active Channels & 1--8; 4 & Adds $8aV/10$, truncates and clamps; depth $a$ defaults to 0. \\
Frequency & $\pm2.5$ oct.; 0 & V/OCT adds octaves; active count also changes actual pitch. \\
FM depth & $-1$--$+1$; 0 & Adds $aV/5$ octaves; unpatched first input is 5\,V. \\
Morph & 1--5; 1 & Adds $5aV/7$ wave positions; morph depth defaults to 0. \\
Level & 0--15; 15 & Multiplies by $V/10$, rounds and clamps; first input normals to 10\,V. \\
\bottomrule
\end{tabularx}

At full active-count depth, 1.25\,V changes the count by one before clipping.
At full morph depth, 1.4\,V moves one waveform position. Knobs supply offsets;
the two depth controls scale and invert their respective CVs.

Pitch, FM and level inputs normal forward along the voice columns.
An inserted cable replaces that signal for the remainder of its chain.
Pitch and FM run at audio rate; level, morph and active-count control update
every 16 host samples. Chip frequency registers limit the tuning range and
make the highest notes coarse.
